%%
%% Copyright (c) 2018-2020 Weitian LI <wt@liwt.net>
%% CC BY 4.0 License
%%
%% Created: 2026.10.22
%%

% Chinese version
\documentclass[zh]{resume}

% File information shown at the footer of the last page
\fileinfo{}

\name{昱蒙}{付}

% \tagline{\texorpdfstring{\icon{\faBinoculars} }{}<position-to-look-for>}
% \tagline{<current-position>}

% \photo[
%   shape=<circular|square>,    % default is circular
%   position=<left|right>,      % default is left
% ]{<width>}{<filename>}
% Example:
% \photo[shape=square]{5.5em}{photo}

\profile{
  \mobile{13903694145}
  \email{13903694145@163.com}
  \university{哈尔滨工业大学}
  \degree{电子信息-计算机技术 \textbullet 硕士}
  \iconlink{\faGithub}{https://github.com/eternal-yangtze/research_claw}{GitHub: research_claw}
  % Custom information:
  % \icontext{<icon>}{<text>}
  % \iconlink{<icon>}{<link>}{<text>}
}

\begin{document}
\makeheader

%======================================================================
% Summary & Objectives
%======================================================================
\begin{abstract}
求职方向：AI Agent、大模型应用开发与 RAG 工程化实习；一周内到岗，每周 5 天，至少 3 个月。哈尔滨工业大学计算机技术硕士在读，预计 2028 年 3 月毕业。独立开发 ResearchOps Agent Runtime，重点解决长任务中的上下文压缩、工具调用链路、任务状态续做、检索证据组织和执行复盘问题；已在其上构建 Coding Agent、RAG 证据检索、Web 多轮问答与离线评测矩阵。
\end{abstract}

%======================================================================
\sectionTitle{专业技能}{\faWrench}
%======================================================================
\begin{itemize}
\item \textbf{Agent 架构}: Python，Agent Runtime/Harness、MicroAgent、多 Agent、LangGraph、Tool Calling、Skills、MCP、任务规划、Memory/上下文管理、Checkpoint 恢复与 Trace 可观测性。
\item \textbf{RAG 与 Eval}: 文档解析与分块、BM25/Embedding 多路召回、Qdrant、Cross-Encoder 重排、MMR、引用校验；LangChain、RAGAS、Recall/nDCG 评测。
\item \textbf{后端与工程}: FastAPI、SQLite、Docker、Linux、Git、pytest；具备沙箱执行、权限边界、状态持久化及故障恢复实践。
\item \textbf{开发基础}: Java、Spring Boot 基础；熟悉 REST API 与前后端协作。
\end{itemize}

%======================================================================
\sectionTitle{教育背景}{\faGraduationCap}
%======================================================================

\begin{itemize}

  \item \textbf{哈尔滨工业大学}
        \quad 计算学部
        \quad 计算机技术
        \quad 硕士
        \hfill 2025.08 -- 2028.06（预计）

  \item \textbf{哈尔滨工业大学}
        \quad 计算学部
        \quad 数据科学与大数据技术
        \quad 学士
        \hfill 2021.08 -- 2025.08


\end{itemize}

%======================================================================
\sectionTitle{项目经历}{\faCode}
%======================================================================
\begin{itemize}
\item \textbf{ResearchOps Agent Runtime（面向科研任务的 Agent 执行平台与 Coding Agent）}

\textbf{项目概述}: 独立设计科研 Agent Runtime，以确定性 Harness 管理模型、工具、状态与工件，以 MicroAgent 承载局部自主推理，贯通“文献取证 $\rightarrow$ 方案规划 $\rightarrow$ 代码生成 $\rightarrow$ 沙箱验证”，并提供 CLI 与 Web 资料问答入口。
\begin{itemize}
\item \textbf{Agent Runtime 与多 Agent 编排}: 设计“外层 Harness + 内层 MicroAgent”架构，外层统一预算、ToolPolicy、状态恢复与 Trace，内层隔离 Prompt、Tools 和反思循环；以 LangGraph 编排 planner、compile gate、critic、final review，通过条件边形成“提出--编译--评审--修订”闭环。
\item \textbf{Agentic RAG 与 MCP}: 以 SQLite/BM25 和 Qdrant/Embedding 多路召回，结合 Fusion、Cross-Encoder 与 MMR 重排；基于官方 MCP SDK 实现 stdio 只读论文服务，Agent 通过 Tool Calling 自主检索、回读原文并反思修订。94 篇论文、12 问上 nDCG@10 由 0.690 提升至 0.749，事实支持率由 75.51\% 提升至 96.15\%（LLM 银标）。
\item \textbf{Skills 与上下文工程}: 实现 \texttt{search\_skills $\rightarrow$ load\_skill $\rightarrow$ invoke\_skill} 三段式按需加载，结合 L0--L3 文献分层摘要、工具观察截断和原文按需回读；同任务对照中初始 Prompt Token 减少 68.7\%，平均总 Token 减少 57.4\%。
\item \textbf{Coding Agent 与沙箱闭环}: 将 LangGraph 研究合同交给 Coding Agent，按模块生成代码，在断网、非 root Docker 沙箱中通过只读挂载和外部 verifier 获取反馈并自动修订；联合验收交付 6/6 模块，Docker 检查 10/10、生成单测 6/6，完成自然语言到可执行源码的闭环。
\item \textbf{分层持久化与故障恢复}: 以 SQLite 存原文、Qdrant 存向量、SQLite Checkpointer 存图状态、JSON 存角色状态；API/网络异常时保留未知费用预留并支持跨进程 Checkpoint 续接，避免重复调用与计费；修复下游费用累加引发的上游工件 Hash 漂移并补充回归覆盖。
\end{itemize}
\end{itemize}

\end{document}
